\section{检验规约与测量结果的复核}
\label{app:gates}

第~\ref{sec:olmo}~节给出了本版本记录的数值结果。以下规约说明，在什么条件下，一次测量可以支持相应的论断。每个比较条件均由一个模型检查点、一个分词器和一个固定接口确定。该版本的完整定义见\src{observer/gate_contract_ru.md}{第二系列实验规约}。

\paragraph{测试集划分与预算。}
将八位输入写作 $\mathbf{x}=LwR$。在 ASCII 字符串 \verb!observer-v2|! 之后拼接六位中间串 $w\in\{0,1\}^6$，保留前导零，再按所得 SHA-256 摘要升序排列。从一开始编号，将第 1--8 名用于开发集（dev），第 9--20 名用于 G1、G2 联合检验，第 21--32 名用于 G3，第 33--44 名用于 G4，第 45--64 名留作备用。对于每个中间串，将四种端点组合与四种历史全部交叉，按字典序呈现 $(L,R,h,e)\in\{0,1\}^4$。每个区块有 16 个条件，开发集共有 128 个条件，上述每个含 12 个区块的留出集有 192 个条件。在开发集上选择措辞、示例、干预层和干预幅度。示例串与评测串分开。首次启用共同的留出集前，先固定待比较的实现。依照已经启用的测试集结果选出的方法，须用新的测试集评估。

\paragraph{取得两种历史。}
首先在温度 1 下，测量单词元续写 \texttt{[space]A}、\texttt{[space]B} 在完整词表上的概率，其中 \texttt{[space]} 表示一个前导空格。准入条件为 $P(A),P(B)\ge 0.01$。使用 \texttt{random.Random(110)} 从条件二元分布 $P(A\mid A\text{ 或 }B)=P(A)/(P(A)+P(B))$ 中有放回地抽样，恰好抽取 1024 次。记录两项原始概率、完整抽样序列及 A 和 B 首次出现的位置。在每条历史中，模型处理实际抽到的第一个对应类型的词元。在处理完该词元之后、反馈到来之前取快照。任一概率低于下限或任一字母未出现，该前缀即记为 \texttt{NOT\_EVALUABLE}。强行插入字母或增加一轮抽样都会改变实验条件。

\paragraph{通用判据。}
对每个必测的数字输出条件，包括中性、恢复和替换对照，都要记录目标、全词表中概率最大的词元、$\ell_1-\ell_0$ 以及 $q=P(0)+P(1)$。用 FP32 或更高精度在全词表上归一化后计算概率。每个条件均须满足数值有限且 $q\ge 0.80$。当差值严格为正或严格为负时，获胜比特的概率 $>0.40$，两者之外任一词元的概率均不超过 $0.20$。差值为零时，该行计入分母，结果计为失败。整块通过要求 16/16 个输出均正确。开发集的通过门槛为 6/8 个整块，留出集为 9/12 个整块。对每个 $e$ 分别公布 $L\ne R$ 时成对求解的结果。这里的 75\% 指整块的通过比例。重复使用同一前缀不会产生彼此独立的模型实例样本。

\paragraph{G0 与 G0R。}
G0 核对权重与分词器的哈希值，检查 A/B/X 和 0/1 在实际上下文中是否各为单个词元，并核对词元 ID 的精确拼接、形状、数据类型和 RoPE 位置。诊断使用一次性克隆。读出过去的选择后，在答案写回前停止。有无该诊断，同一主查询都必须给出完全相同的全词表 logits。两条分支的原始快照保留所有有效 K/V 的字节、偏移量与可见位置顺序，包括环形缓冲区的索引。比较对象是有效位置对应的逻辑数据。此范围以外分配器的闲置容量不计入检查。

在进行筛选之前，G0R 对所有层的选择位置执行可逆替换，此项检验独立于行为评分。若共同前史为 $U$，两份快照可写为 $C_A=(U,a)$、$C_B=(U,b)$。对每种反馈文本和每个新输入串，检验路径 $A\to B\to A$ 以及从 B 出发的镜像路径。在正向替换所得快照上执行反向替换，读出操作使用其克隆。在中点与未经修改的供体历史比较，在终点与起始历史比较。以相同布局核对元数据、有效 K/V 和 logits。形状须一致，数值须有限，原始数据类型下的字节须逐一相同。将 logits 扩展为 FP32 后，最大绝对差须为零。检验前后均核对参考快照。首次 Base 实验固定使用原始 \texttt{prefix(8)}。替换前验证共同前缀相同，并验证选择位置至少有一个有效字节不同。若两条分支在该位置已无法区分，该前缀的 G0R 结果为 \texttt{NOT\_EVALUABLE}。开发集有 128 行，每行比较两次 logits，共 256 次检查。

\paragraph{G1 与 G2。}
G1 在未经修改的历史上检验利用文本反馈进行的计算。除必测的数字输出探测外，还用彼此独立的副本执行 D1（在反馈前复述 $h$）、D2（不提供新输入串，分别对四组 $h,e$ 给出修正后的规则）和 D3（在新回合中应用明确指定的规则）。字母输出的诊断结果与主数字任务分别报告。文本对照的通过条件由输出格式与整块通过比例共同确定。对 G2，在 G0、G0R 和 G1 之后，在第 9--20 名组成的测试集上重复供体比较和恢复比较，并添加一个共同缓存，将等长的中性词元 X 放在同一位置。由 $P_B(U,a)=(U,b)$ 可知替换结果与未经修改的供体相同。因此技术检验通过时的判定为 \texttt{REPLAY\_OK}。未经修改的历史上的行为效应已由 G1 评估。

\paragraph{G3。}
通过 G1 并取得 \texttt{REPLAY\_OK} 后，用一次激活干预取代文本信号。干预发生在中性词元的 K/V 写入之前。测试前固定词元、位置、层、干预点、向量来源和幅度。在执行运算的数据类型中表示一个有限非零向量 $v$。通过精确变号得到脉冲 $v_0=v$、$v_1=-v$。对于相同的 $e$，两条历史收到逐字节相同的张量。检查形状、类型、逐元素互为相反数，并在 FP64 下按同一运算顺序验证两者的欧几里得范数非零且相等。另对每条分支记录实际舍入后的增量 $\operatorname{round}(r+v_e)-r$ 的范数。一次加法后即关闭干预。干预为零时，在固定 $h,\mathbf{x}$ 下，两条 $e$ 分支必须具有相同的输入和 logits，并公布其准确率及概率质量。要判定脉冲实验通过，须符合通用格式要求，且在 G3 测试集上至少解出 9/12 个整块。

\paragraph{G4。}
在呈现 $\mathbf{x}$ 之前，四个供体在 $h,e$ 之后处理相同的中性序列。仅将此序列中预先指定位置和层的 K/V 转移至新模板，记作 $M_{he}$。四个接收方的其他有效位置、词元、长度和元数据均相同。接收方无法直接访问 $h,e$ 或供体历史的其他后继状态。中性状态 $M_*$ 由共同的中性历史一次确定。文本通道以 G1、G2 为前提，脉冲通道另须通过 G3。对 12 个中间串，以下条件须同时满足：
\begin{enumerate}
\item 给接收方完整历史及相同的中间序列，它须满足格式要求并至少解出 9/12 个整块。若该对照失败，转移未成功仍不足以判定供体中缺少所需信息。
\item 插入 $M_*$ 后，在固定 $\mathbf{x}$ 下，四个输入、有效 K/V 和 logits 须逐字节相同。在同一个派生快照上，先插入其他供体的状态，再换回自身的 $M_{he}$，须恢复原始 logits。
\item 转移各自的 $M_{he}$ 时，须满足格式要求及 9/12 个整块的门槛。每一块覆盖全部四个供体与四种端点组合。
\item 对 $L\ne R$ 的 96 个条件，分别以 $C_M$ 和 $C_*$ 表示转移和擦除时的正确答案数。要求格式有效，$C_*\le48$，擦除后整块通过数为 0/12，且 $(C_M-C_*)/96\ge 0.25$。分别统计由错变对和由对变错的条件。若每行必须给出一个严格的二元答案，共同输入的正确数恰为 48/96。转移时解出九个整块，保证在端点不同的条件上至少答对 72/96。
\end{enumerate}

\paragraph{G4-1。}
在看到任何 $\mathbf{x}$ 之前，发送方通过全词表贪心解码自由选择一个词元。只接受预先指定的两个单词元结尾 A/B。若最高分并列或胜出词元不在 A/B 中，记作 \texttt{FAIL\_FORMAT}，不再重试。按 $00,01,10,11$ 的历史顺序记录四条消息。原始 Base 规约固定了 \texttt{prefix(8)}、字面反馈 correct/wrong 和记录修正规则的请求。实际测量的两步 Instruct 条件有独立的提示词和附加门槛 $P(A)+P(B)\ge0.80$，详见\src{observer/two-step-interface-04.json}{方案 A 的清单}。分别处理两个模型检查点的条件。

复现原始 Base 规约时，在发送方选定后传入 \verb!\nFeedback: correct! 或 \verb!\nFeedback: wrong!，随后传入 \verb!\nRecord the corrected rule. Rule:!。其中 \verb!\n! 表示一个换行符。新接收方的前缀为以下两行：
\begin{verbatim}
Rule A selects the first bit. Rule B selects the last bit.
Rule:
\end{verbatim}
冒号后立即插入带前导空格的所选词元 ID。接着是 \verb!\nBits: !、用单个空格分隔的八个比特，以及末尾有一个空格的 \verb!\nAnswer: !。

在新回合中，接收方仅获得选定的 ID，并重新计算其 K/V。将共同前缀及查询串的词元 ID 与完整文本的分词结果核对。外部接口不传递发送方的 K/V、logits、分支编号、响应时间或抽样参数。禁止依据 $h,e,p$、正确答案或其替代指标来选择载体或分支。擦除时，在所有世界中都使用相同输入，并传入同一个预先固定的词元 A。判定为 \texttt{ZERO\_ERROR\_ON\_PANEL} 须有四条合规消息，通过通道边界与擦除检查，满足数字输出格式，且解出 12/12 个整块。若联合检验 G4 和 G4-1，两种方案都须在首次启用第 33--44 名组成的测试集之前固定。

\paragraph{状态优先级与可复现性。}
出现技术错误、非有限数值或违反所要求的逐字节同一性，记为 \texttt{INVALID}。其后按优先级递减依次为 \texttt{NOT\_EVALUABLE}、\texttt{FAIL\_FORMAT} 和 \texttt{FAIL\_DEV}/\texttt{FAIL\_TEST}。因前序条件未满足而未执行的检查记为 \texttt{NOT\_RUN}。修改接口或更换模型检查点，即构成新的实验条件。

从固定版本的软件包根目录运行下列命令，二者各有用途。
\begin{verbatim}
python3 -B observer/verify.py
python3 -B observer/verify_part_iii.py
\end{verbatim}
第一个验证程序重新计算有限构造，检查历史语料发布版的清单、登记表各项的地址，以及能否从 SVG 恢复表格。第二个验证程序读取已保存的开发集报告，检查各条件组合及记录的概率和 logits，重新算出准确率、完整区块数、相反历史配对的结果、诊断结果与擦除效应。对 G0/G0R，它再次核对元数据以及成对比较原始类型 logits 时记录的哈希值。重新运行 GPU 并对新算出的张量逐字节比较，将构成一次新的测量。
